%!TEX root = ../../main.tex
\section{존재 게이트 $R$, $B$, $M$}
\label{sec:eng-gate}

킬만으로는 양 팀이 그 직전에 싸울 위치에 있었는지 알 수 없다. 에피소드는 첫 킬 $B$초 전에 둔 예측 시점 $\tcut$에서 양 팀 각각 살아 있는 챔피언 $M = \constM$명 이상이 첫 킬 위치(\emph{앵커})로부터 $R$ 안에 있을 때에만 \emph{교전}이 된다. 생존은 $\tcut$ 이전의 마지막 1분 프레임에서 읽고, 위치는 1분 프레임 사이를 보간한 5초 격자에서 읽되 킬 참가자의 경로는 그 킬 위치 쪽으로 끌어당긴다. Ke 등 \cite{ke2022}은 한쪽에 두 명 이상을 요구하지만 본 논문은 양쪽 모두에 요구한다.

$R$과 $B$는 위치가 1분마다 기록되는 Match-V5에서 추정할 수 없다. 따라서 둘을 게임 자체가 밀접한 목적에 쓰는 값에 가정으로 고정한다. $R = \constR$ 단위는 챔피언 사망 시 상대 챔피언들이 경험치를 나누어 갖는 반경이고, $B = \constB$초는 소환사의 협곡에서 어시스트가 인정되는 창이다 \cite{lolwiki2026experience,lolwiki2026assist}. $B$는 Tot 등 \cite{tot2021camera}이 쓴 OpenDota 처리기가 첫 사망 전에 전투 창을 여는 오프셋과도 같다 \cite{opendota2018processteamfights}. 패치 25.14--25.16의 공식 노트와 핫픽스에는 경험치 공유 반경이나 어시스트 창의 변경이 없다. ``변경이 기재되지 않음''이 이 검사의 강도이다. 패치 노트는 변경을 빠뜨릴 수 있으므로 값이 불변임이 검증된 것은 아니다. $B$는 전투의 시작이 아니라 앞선 여유(lead)이다.

$R$과 $B$를 가정으로 고정하였으므로 둘이 코퍼스를 얼마나 움직이는지 보고한다. 2{,}971 경기의 81{,}914 킬 에피소드에서 교전으로 인정되는 비율은 (\constR{} 단위, \constB{}초)에서 \gateShareAdopted{}\%, 선행 연구의 점 (1{,}800 단위, 10초)에서 \gateShareCog{}\%이다. $R \in \{1{,}350, 1{,}600, 1{,}800, 2{,}500\}$과 $B \in \{5, 10, 15\}$에 걸쳐 그 비율은 6.5\%에서 41.1\%까지 여섯 배 변한다. 게이트는 교전 집합의 크기를 결정하고 예측 가능성도 움직인다. 선행 연구 계보의 교전 승자 라벨로 게이트만 바꿔 21만 경기를 다시 돌리면, (1{,}800 단위, 10초)에서 코퍼스는 \numEngGateCog{} 교전으로 1.8배가 되고 두 게이트가 모두 받아들인 교전에서의 짝 AUC 차이는 $+0.0039$였으며, $M = 3$에서는 \numEngMthree{} 교전으로 줄고 공통 교전에서 $-0.0066$이었다. 즉 더 높은 절대 AUC는 어떤 교전이 남는가에서 온다. 예측 결과는 게이트에 불변이 아니며, $R$과 $B$는 그것이 주는 AUC가 아니라 위의 게임 규칙에 의해, $M$은 관례에 의해 고정된 채로 둔다.

앵커는 에피소드 킬들의 무게중심이 아니라 첫 킬의 위치이다. 선행 연구 상수 아래 파일럿에서 앵커를 무게중심으로 옮기면 6.0\%의 에피소드에서 게이트 판정이 바뀌었고, 이 실험은 v3.3 상수로 반복되지 않았다.
